\section{Where trajectory length changes the trade-off}
\label{sec:horizon}

The aggregate result hides substantial differences across task trajectories.
We use the control-defined cohorts from Section~\ref{sec:setup}: short
($T\leq8$), medium ($8<T\leq20$), and long ($T>20$), where $T$ is each
task's median full-history trajectory length. Assignment is held fixed across
memory policies. A policy's own trajectory length is an outcome, not the
variable used to select its cohort.

\begin{table}[htbp]
\centering\small
\caption{Control-defined trajectory-length cohorts: quality and cost.
All retained outcomes are included. Cost ratio is the median same-task,
same-repeat ratio to full history.}
\label{tab:horizon-main}
\resizebox{\textwidth}{!}{\input{tables/new-horizon}}
\end{table}

State-only success is lower in every cohort. The largest deficit occurs in
the long cohort: \HorizonLongModelStateRate\% against
\HorizonLongModelControlRate\% for full history. Retaining $k=6$ restores
\HorizonLongModelHistorySixRate\% (Table~\ref{tab:horizon-main}). Its cost ratio
is \HorizonLongModelHistorySixAllCostRatio{}, compared with
\HorizonLongModelStateAllCostRatio{} for state alone. Cost per observed success
is respectively \$\HorizonLongModelHistorySixCostPerSuccess{} and
\$\HorizonLongModelStateCostPerSuccess{}, versus
\$\HorizonLongModelControlCostPerSuccess{} for the control.

\begin{table}[htbp]
\centering\small
\caption{The same trajectory-length cohorts: median context and token
measurements. Cache hit is the median per-execution fraction $K/(U+K)$;
token columns are separate medians.}
\label{tab:horizon-cache}
\resizebox{\textwidth}{!}{\input{tables/new-horizon-resources}}
\end{table}

Long trajectories also strengthen the full-history cache advantage.
Its median cache-hit rate rises from \HorizonShortModelControlCacheHit\%
in the short cohort to \HorizonLongModelControlCacheHit\% in the long cohort;
state alone and $k=6$ reach \HorizonLongModelStateCacheHit\% and
\HorizonLongModelHistorySixCacheHit\% there. Cached-token volume alone cannot
show this difference: hit rate also depends on uncached input.

These descriptive cohorts identify where a history window is useful, not a
universal switching threshold. Short-cohort tasks can have long outlier
repeats; workload-level efficiency and median matched costs therefore need
not rank policies alike. Trajectories are policy-dependent, and no controlled
intervention on task horizon establishes a causal crossover here.
